\section{质心力学}

\subsection{质心动量定理}

\subsubsection{质心}

\begin{definition}
    对于一组质点，质量分别为 $m_1, m_2,\dots, m_n$，位移分别为 $\vect{r}_1, \vect{r}_2, \dots, \vect{r}_n$，则定义其\textbf{质心}为一个质量为
    $$
    M = \sum_{i=1}^n m_i
    $$
    位移为
    $$
    \vec{r}_c = \frac{1}{M} \sum_{i=1}^n m_i \vect{r}_i
    $$
    的虚拟的质点。
\end{definition}

\begin{definition}
    质心的\textbf{动量}定义为质点组中质点的动量之和，即
    $$
    \vect{p}_c = \sum_{i=1}^n \vect{p}_i = \sum_{i=1}^n m_i \vect{v}_i = M \vect{v}_c
    $$
\end{definition}

根据上述定义，可以得到质心的动量定理：

\begin{theorem}
    质心的动量变化量等于质点组所受合外力的冲量。
\end{theorem}

\subsubsection{质心参考系}

当我们称原点固定在质心上的平动参考系为\textbf{质心}参考系。

\subsection{质心动能定理}

根据
$$
\begin{aligned}
    E_k & = \sum_{i=1}^n \frac{1}{2} m_i \vect{v}_i^2 \\
    & = \sum_{i=1}^n \frac{1}{2} m_i \vect{v}_c^2 + \sum_{i=1}^n \frac{1}{2} m_i \vect{v}_i'^2 + \sum_{i=1}^n m_i \vect{v}_c \cdot \vect{v}_i' \\
    & = \frac{1}{2} M \vect{v}_c^2 + \sum_{i=1}^n \frac{1}{2} m_i \vect{v}_i'^2 + \sum_{i=1}^n m_i \vect{v}_c \cdot \vect{v}_i'
\end{aligned}
$$
其中
$$
\begin{aligned}
    \sum_{i=1}^n m_i \vect{v}_c \cdot \vect{v}_i' = \vect{v}_c \cdot \sum_{i=1}^n m_i \vect{v}_i' = \vect{v}_c \cdot \vect{p}_c' = 0
\end{aligned}
$$
因此：
$$
E_k = \frac{1}{2} M \vect{v}_c^2 + \sum_{i=1}^n \frac{1}{2} m_i \vect{v}_i'^2
$$
于是我们得到：

\begin{theorem}[质心动能定理]
    质点组的动能等于质心动能与质点组相对于质心参考系的动能之和。
\end{theorem}

\subsection{质心角动量定理}

\subsubsection{质心的角动量}

\begin{definition}
    定义质心的角动量为
    $$
    \vect{L}_c = \vect{r}_c \times \vect{p}_c
    $$
\end{definition}

进行一些推导：
$$
\begin{aligned}
    \vect{L} & = \sum_{i = 1}^n m_i \vect{r}_i \times \vect{v}_i \\
    & = \sum_{i = 1}^{n} m_i (\vect{r}_c + \vect{r}_i') \times (\vect{v}_c + \vect{v}_i') \\
    & = \sum_{i = 1}^{n} m_i \vect{r}_c \times \vect{v}_c + \sum_{i = 1}^{n} m_i \vect{r}_c \times \vect{v}_i' + \sum_{i = 1}^{n} m_i \vect{r}_i' \times \vect{v}_c + \sum_{i = 1}^{n} m_i \vect{r}_i' \times \vect{v}_i'
\end{aligned}
$$
中间两项均为 $\vect{0}$，化简得到：
$$
\vect{L} = \vect{L}_c + \sum_{i = 1}^{n} \vect{L}_i'
$$

\subsection{有心运动方程和约化质量}

